\documentclass[twocolumn]{article}
\usepackage[utf8]{inputenc}
\usepackage{geometry}
\usepackage{amsmath}
\usepackage{hyperref}

\geometry{a4paper, margin=1in}

\title{\textbf{Distributed SNP: Executive Summary}}
\author{}
\date{}

\begin{document}

\maketitle

\section*{Project Overview}
The \textbf{Distributed SNP} project is a high-performance C++ library designed for the simulation of Spiking Neural P (SNP) Systems. By reformulating the non-deterministic evolution of SNP systems into linear algebra operations, the library leverages modern parallel computing architectures to achieve scalable simulations. The core state evolution is governed by the matrix equation:
\[ C(k+1) = C(k) + Sp(k) \times M \]
where $C$ represents the configuration vector, $Sp$ the spiking vector, and $M$ the transition matrix.

\section*{Key Features}
\begin{itemize}
    \item \textbf{Hybrid Parallelism}: The system utilizes a hybrid parallel approach, combining \textbf{MPI} (Message Passing Interface) for distributed memory parallelism across nodes and \textbf{CUDA} for shared memory parallelism on GPUs.
    \item \textbf{Versatile Backends}: The library supports multiple simulation backends to cater to different hardware configurations:
    \begin{itemize}
        \item \textbf{CPU}: OpenMP-accelerated implementation for standard multi-core processors.
        \item \textbf{Naive CUDA}: Single-GPU implementation for dense systems.
        \item \textbf{Distributed MPI+CUDA}: Multi-GPU implementation scaling across networked nodes.
        \item \textbf{Sparse CUDA}: Optimized implementation for systems with sparse connectivity matrices.
    \end{itemize}
\end{itemize}

\section*{Architecture}
The codebase follows a modular architecture designed for extensibility and maintainability:
\begin{itemize}
    \item \textbf{Interface-Driven Design}: Core functionalities are abstracted behind interfaces such as \texttt{ISnpSimulator} for simulation logic and \texttt{IMatrixOps} for linear algebra operations. This allows for seamless swapping of backends without altering the application logic.
    \item \textbf{Factory Pattern}: Concrete implementations are instantiated via factory functions, ensuring loose coupling between components.
    \item \textbf{Separation of Concerns}: The simulation rules (SNP logic) are strictly separated from the computational engine (Matrix Operations), enabling independent optimization of the math kernel.
\end{itemize}

\section*{Codebase Structure}
\begin{itemize}
    \item \texttt{src/snp}: Contains the core SNP simulation logic, including the simulator interfaces and concrete implementations (e.g., \texttt{CudaMpiSnpSimulator}, \texttt{SparseCudaSnpSimulator}), as well as graph partitioners (Linear, Louvain).
    \item \texttt{src/linear\_algebra}: Houses the matrix operation engines. Implementations include \texttt{CpuMatrixOps}, \texttt{CudaMatrixOps}, and \texttt{MpiCudaMatrixOps}.
    \item \texttt{src/sort}: Demonstrates the library's capabilities through an application of SNP systems for sorting algorithms (Bitonic sort).
\end{itemize}

\section*{Build \& Test Workflow}
\begin{itemize}
    \item \textbf{Build System}: The project uses \textbf{CMake} (3.18+) for cross-platform build configuration.
    \item \textbf{Testing}: A comprehensive test suite is implemented using \textbf{Google Test}, covering unit tests for matrix operations and integration tests for simulators.
    \item \textbf{Benchmarking}: Performance is evaluated using a dedicated benchmarking suite (Google Benchmark) and visualization scripts located in the \texttt{benchmark/} and \texttt{scripts/} directories.
\end{itemize}

\end{document}
